\chapter{Nerve Pain (Neuropathic Pain)}

\section*{Definition}
Pain caused by damage or dysfunction of the somatosensory nervous system, characterized by burning, shooting, electric shock-like, or stabbing sensations.

\section*{Pathophysiology}
Nerve injury leads to abnormal ectopic discharge from damaged axons, altered ion channel expression, central sensitization in the spinal cord dorsal horn, and reduced descending inhibitory control. Neuroplastic changes in the brain amplify pain perception, making neuropathic pain often chronic and difficult to treat.

\section*{Etiology}
\begin{itemize}
  \item Diabetic peripheral neuropathy
  \item Postherpetic neuralgia (shingles)
  \item Trigeminal neuralgia
  \item Sciatica or lumbar radiculopathy
  \item Carpal tunnel syndrome
  \item Multiple sclerosis
  \item Chemotherapy-induced peripheral neuropathy
  \item HIV-associated neuropathy
  \item Alcohol-related neuropathy
  \item Complex regional pain syndrome (CRPS)
\end{itemize}

\section*{Indications for Evaluation}
Seek care if:
\begin{itemize}
  \item Severe or worsening symptoms
  \item New or progressive burning or shooting pain, pain with weakness, numbness, or changes in bladder/bowel function
  \item Accompanied by other concerning signs
\end{itemize}

\section*{Related Symptoms}
\begin{itemize}
  \item Numbness
  \item Tingling
  \item Burning sensation
  \item Weakness
  \item Allodynia (pain to light touch)
\end{itemize}
\textbf{Note:} Always consider serious underlying conditions when evaluating persistent nerve pain (neuropathic pain).

\section*{Self-Care / Home Management}
\begin{itemize}
  \item Rest and avoid activities that may aggravate the condition.
  \item Maintain adequate hydration and a balanced diet.
  \item Over-the-counter remedies may provide symptomatic relief (consult a pharmacist or doctor).
  \item Monitor symptoms and seek professional advice if they persist or worsen.
  \item Avoid self-medicating with prescription medications without medical guidance.
\end{itemize}

\section*{Diagnostic Approach}
\begin{itemize}
  \item A thorough history and physical examination are the first steps in evaluation.
  \item Laboratory tests (blood work, urinalysis) may be ordered based on clinical suspicion.
  \item Imaging studies (X-ray, ultrasound, CT, MRI) may be indicated when structural pathology is suspected.
  \item Specialist referral is arranged when the diagnosis remains uncertain or requires specific expertise.
\end{itemize}

\section*{Treatment / Management Overview}
\begin{itemize}
  \item Treatment targets the underlying cause identified through diagnostic evaluation.
  \item Supportive care and symptom management are provided while awaiting definitive therapy.
  \item Medications, lifestyle modifications, and physical therapies may be prescribed as appropriate.
  \item Follow-up ensures treatment effectiveness and allows timely adjustments to the care plan.
\end{itemize}

\clearpage